%% Generated by Sphinx.
\def\sphinxdocclass{report}
\documentclass[letterpaper,10pt,english]{sphinxmanual}
\ifdefined\pdfpxdimen
   \let\sphinxpxdimen\pdfpxdimen\else\newdimen\sphinxpxdimen
\fi \sphinxpxdimen=.75bp\relax
\ifdefined\pdfimageresolution
    \pdfimageresolution= \numexpr \dimexpr1in\relax/\sphinxpxdimen\relax
\fi
%% let collapsible pdf bookmarks panel have high depth per default
\PassOptionsToPackage{bookmarksdepth=5}{hyperref}

\PassOptionsToPackage{booktabs}{sphinx}
\PassOptionsToPackage{colorrows}{sphinx}

\PassOptionsToPackage{warn}{textcomp}
\usepackage[utf8]{inputenc}
\ifdefined\DeclareUnicodeCharacter
% support both utf8 and utf8x syntaxes
  \ifdefined\DeclareUnicodeCharacterAsOptional
    \def\sphinxDUC#1{\DeclareUnicodeCharacter{"#1}}
  \else
    \let\sphinxDUC\DeclareUnicodeCharacter
  \fi
  \sphinxDUC{00A0}{\nobreakspace}
  \sphinxDUC{2500}{\sphinxunichar{2500}}
  \sphinxDUC{2502}{\sphinxunichar{2502}}
  \sphinxDUC{2514}{\sphinxunichar{2514}}
  \sphinxDUC{251C}{\sphinxunichar{251C}}
  \sphinxDUC{2572}{\textbackslash}
\fi
\usepackage{cmap}
\usepackage[T1]{fontenc}
\usepackage{amsmath,amssymb,amstext}
\usepackage{babel}



\usepackage{tgtermes}
\usepackage{tgheros}
\renewcommand{\ttdefault}{txtt}



\usepackage[Bjarne]{fncychap}
\usepackage[,numfigreset=2,mathnumfig,mathnumsep={.}]{sphinx}
\sphinxsetup{
verbatimwithframe=false,
VerbatimColor={rgb}{0.96,0.96,0.96},
VerbatimBorderColor={rgb}{1,1,1}
}
\fvset{fontsize=auto}
\usepackage{geometry}


% Include hyperref last.
\usepackage{hyperref}
% Fix anchor placement for figures with captions.
\usepackage{hypcap}% it must be loaded after hyperref.
% Set up styles of URL: it should be placed after hyperref.
\urlstyle{same}


\usepackage{sphinxmessages}




\title{Math 335 Homework 1 Solutions}
\date{Jun 12, 2026}
\release{0.1}
\author{Biraj Borgohain}
\newcommand{\sphinxlogo}{\vbox{}}
\renewcommand{\releasename}{Release}
\makeindex
\begin{document}

\ifdefined\shorthandoff
  \ifnum\catcode`\=\string=\active\shorthandoff{=}\fi
  \ifnum\catcode`\"=\active\shorthandoff{"}\fi
\fi

\pagestyle{empty}
\sphinxmaketitle
\pagestyle{plain}

\pagestyle{normal}
\phantomsection\label{\detokenize{Language/Science_language/ODE_calculus/course_MATH_3035/homework_assignments/index_pdf::doc}}



\section{Submited Homework Assignment MATH335 HW1}
\label{\detokenize{Language/Science_language/ODE_calculus/course_MATH_3035/homework_assignments/index_pdf:submited-homework-assignment-math335-hw1}}

\subsection{By Biraj Borgohain (900369162)}
\label{\detokenize{Language/Science_language/ODE_calculus/course_MATH_3035/homework_assignments/index_pdf:by-biraj-borgohain-900369162}}

\section{Solution of Problem 1:}
\label{\detokenize{Language/Science_language/ODE_calculus/course_MATH_3035/homework_assignments/index_pdf:solution-of-problem-1}}

\begin{savenotes}\sphinxattablestart
\sphinxthistablewithglobalstyle
\centering
\begin{tabular}[t]{\X{10}{70}\X{10}{70}\X{15}{70}\X{35}{70}}
\sphinxtoprule
\sphinxstyletheadfamily 
\sphinxAtStartPar
Equation
&\sphinxstyletheadfamily 
\sphinxAtStartPar
Order
&\sphinxstyletheadfamily 
\sphinxAtStartPar
Type
&\sphinxstyletheadfamily 
\sphinxAtStartPar
Reason
\\
\sphinxmidrule
\sphinxtableatstartofbodyhook\begin{enumerate}
\sphinxsetlistlabels{\alph}{enumi}{enumii}{(}{)}%
\item {} 
\end{enumerate}
&
\sphinxAtStartPar
1
&
\sphinxAtStartPar
Nonlinear
&
\sphinxAtStartPar
Contains \(\sin(y)\)
\\
\sphinxhline\begin{enumerate}
\sphinxsetlistlabels{\alph}{enumi}{enumii}{(}{)}%
\setcounter{enumi}{1}
\item {} 
\end{enumerate}
&
\sphinxAtStartPar
1
&
\sphinxAtStartPar
Linear
&
\sphinxAtStartPar
\(y\) and \(y'\) appear linearly
\\
\sphinxhline\begin{enumerate}
\sphinxsetlistlabels{\alph}{enumi}{enumii}{(}{)}%
\setcounter{enumi}{2}
\item {} 
\end{enumerate}
&
\sphinxAtStartPar
3
&
\sphinxAtStartPar
Linear
&
\sphinxAtStartPar
All terms are linear in \(y\) and its derivatives
\\
\sphinxhline\begin{enumerate}
\sphinxsetlistlabels{\alph}{enumi}{enumii}{(}{)}%
\setcounter{enumi}{3}
\item {} 
\end{enumerate}
&
\sphinxAtStartPar
3
&
\sphinxAtStartPar
Nonlinear
&
\sphinxAtStartPar
Contains \(y^4\)
\\
\sphinxhline\begin{enumerate}
\sphinxsetlistlabels{\alph}{enumi}{enumii}{(}{)}%
\setcounter{enumi}{4}
\item {} 
\end{enumerate}
&
\sphinxAtStartPar
1
&
\sphinxAtStartPar
Nonlinear
&
\sphinxAtStartPar
Contains \(1/y\)
\\
\sphinxhline\begin{enumerate}
\sphinxsetlistlabels{\alph}{enumi}{enumii}{(}{)}%
\setcounter{enumi}{5}
\item {} 
\end{enumerate}
&
\sphinxAtStartPar
1
&
\sphinxAtStartPar
Nonlinear
&
\sphinxAtStartPar
Contains \(y^2\)
\\
\sphinxbottomrule
\end{tabular}
\sphinxtableafterendhook\par
\sphinxattableend\end{savenotes}




\section{Solution of Problem 2(a) and (b)}
\label{\detokenize{Language/Science_language/ODE_calculus/course_MATH_3035/homework_assignments/index_pdf:solution-of-problem-2-a-and-b}}

\subsection{Solution 2(a)}
\label{\detokenize{Language/Science_language/ODE_calculus/course_MATH_3035/homework_assignments/index_pdf:solution-2-a}}
\sphinxAtStartPar
Given
\begin{equation*}
\begin{split}y' = y^2 - 1, \qquad y(0)=0\end{split}
\end{equation*}\begin{itemize}
\item {} 
\sphinxAtStartPar
Equilibrium points:
\begin{equation*}
\begin{split}y^2 - 1 = 0 \quad \Rightarrow \quad y=\pm 1\end{split}
\end{equation*}
\item {} 
\sphinxAtStartPar
Sign of \(y'\):
\begin{itemize}
\item {} 
\sphinxAtStartPar
\(y>1 \Rightarrow y'>0\) (solution moves upward)

\item {} 
\sphinxAtStartPar
\(-1<y<1 \Rightarrow y'<0\) (solution moves downward)

\item {} 
\sphinxAtStartPar
\(y<-1 \Rightarrow y'>0\) (solution moves upward)

\end{itemize}

\item {} 
\sphinxAtStartPar
Stability:
\begin{itemize}
\item {} 
\sphinxAtStartPar
\(y=-1\) is a stable equilibrium.

\item {} 
\sphinxAtStartPar
\(y=1\) is an unstable equilibrium.

\end{itemize}

\item {} 
\sphinxAtStartPar
Since \(y(0)=0\) lies in \((-1,1)\), we have \(y'<0\), so the solution decreases toward the stable equilibrium \(y=-1\).

\end{itemize}

\sphinxAtStartPar
Therefore,
\begin{equation*}
\begin{split}\lim_{t\to\infty} y(t) = -1.\end{split}
\end{equation*}
\sphinxAtStartPar
For y’ = y\textasciicircum{}2 \sphinxhyphen{} 1:

\sphinxAtStartPar
Equilibrium points:
\begin{equation*}
\begin{split}y^2 - 1 = 0\end{split}
\end{equation*}\begin{equation*}
\begin{split}y = \pm 1\end{split}
\end{equation*}
\sphinxAtStartPar
Representative slopes:


\begin{savenotes}\sphinxattablestart
\sphinxthistablewithglobalstyle
\centering
\begin{tabulary}{\linewidth}[t]{TT}
\sphinxtoprule
\sphinxstyletheadfamily 
\sphinxAtStartPar
y
&\sphinxstyletheadfamily 
\sphinxAtStartPar
y’
\\
\sphinxmidrule
\sphinxtableatstartofbodyhook
\sphinxAtStartPar
2
&
\sphinxAtStartPar
3
\\
\sphinxhline
\sphinxAtStartPar
1
&
\sphinxAtStartPar
0
\\
\sphinxhline
\sphinxAtStartPar
0
&
\sphinxAtStartPar
\sphinxhyphen{}1
\\
\sphinxhline
\sphinxAtStartPar
\sphinxhyphen{}1
&
\sphinxAtStartPar
0
\\
\sphinxhline
\sphinxAtStartPar
\sphinxhyphen{}2
&
\sphinxAtStartPar
3
\\
\sphinxbottomrule
\end{tabulary}
\sphinxtableafterendhook\par
\sphinxattableend\end{savenotes}

\sphinxAtStartPar
Thus:
\begin{itemize}
\item {} 
\sphinxAtStartPar
y \textgreater{} 1  =\textgreater{} y’ \textgreater{} 0 (solutions move upward)

\item {} 
\sphinxAtStartPar
\sphinxhyphen{}1 \textless{} y \textless{} 1 =\textgreater{} y’ \textless{} 0 (solutions move downward)

\item {} 
\sphinxAtStartPar
y \textless{} \sphinxhyphen{}1 =\textgreater{} y’ \textgreater{} 0 (solutions move upward)

\end{itemize}

\sphinxAtStartPar
Since y(0)=0 lies in (\sphinxhyphen{}1,1), the solution decreases and approaches y=\sphinxhyphen{}1 as t\(\rightarrow\)\(\infty\).


\subsubsection{Slope field plot}
\label{\detokenize{Language/Science_language/ODE_calculus/course_MATH_3035/homework_assignments/index_pdf:slope-field-plot}}
\begin{figure}[htbp]
\centering
\capstart

\noindent\sphinxincludegraphics[width=0.800\linewidth]{{slope_field_q2a}.png}
\caption{Slope field for \(y' = y^2 - 1\).}\label{\detokenize{Language/Science_language/ODE_calculus/course_MATH_3035/homework_assignments/index_pdf:id8}}\end{figure}


\subsection{Solution 2(b)}
\label{\detokenize{Language/Science_language/ODE_calculus/course_MATH_3035/homework_assignments/index_pdf:solution-2-b}}
\sphinxAtStartPar
Given
\begin{equation*}
\begin{split}y' = t^2 - 1, \qquad y(0)=0\end{split}
\end{equation*}
\sphinxAtStartPar
Representative slopes:


\begin{savenotes}\sphinxattablestart
\sphinxthistablewithglobalstyle
\centering
\begin{tabulary}{\linewidth}[t]{TT}
\sphinxtoprule
\sphinxstyletheadfamily 
\sphinxAtStartPar
t
&\sphinxstyletheadfamily 
\sphinxAtStartPar
y’
\\
\sphinxmidrule
\sphinxtableatstartofbodyhook
\sphinxAtStartPar
\sphinxhyphen{}2
&
\sphinxAtStartPar
3
\\
\sphinxhline
\sphinxAtStartPar
\sphinxhyphen{}1
&
\sphinxAtStartPar
0
\\
\sphinxhline
\sphinxAtStartPar
0
&
\sphinxAtStartPar
\sphinxhyphen{}1
\\
\sphinxhline
\sphinxAtStartPar
1
&
\sphinxAtStartPar
0
\\
\sphinxhline
\sphinxAtStartPar
2
&
\sphinxAtStartPar
3
\\
\sphinxbottomrule
\end{tabulary}
\sphinxtableafterendhook\par
\sphinxattableend\end{savenotes}

\sphinxAtStartPar
Sign of \(y'\):
\begin{itemize}
\item {} 
\sphinxAtStartPar
\(t < -1 \Rightarrow y' > 0\) (upward slopes)

\item {} 
\sphinxAtStartPar
\(t = -1 \Rightarrow y' = 0\) (horizontal slopes)

\item {} 
\sphinxAtStartPar
\(-1 < t < 1 \Rightarrow y' < 0\) (downward slopes)

\item {} 
\sphinxAtStartPar
\(t = 1 \Rightarrow y' = 0\) (horizontal slopes)

\item {} 
\sphinxAtStartPar
\(t > 1 \Rightarrow y' > 0\) (upward slopes)

\end{itemize}

\sphinxAtStartPar
Since \(y(0)=0\), the solution initially decreases because
\begin{equation*}
\begin{split}y'(0)=0^2-1=-1.\end{split}
\end{equation*}
\sphinxAtStartPar
For large \(t\),
\begin{equation*}
\begin{split}y' = t^2 - 1 > 0,\end{split}
\end{equation*}
\sphinxAtStartPar
so the solution eventually turns upward and increases without bound.

\sphinxAtStartPar
Therefore,
\begin{equation*}
\begin{split}\lim_{t\to\infty} y(t)=\infty.\end{split}
\end{equation*}

\subsubsection{Slope Field plot}
\label{\detokenize{Language/Science_language/ODE_calculus/course_MATH_3035/homework_assignments/index_pdf:id1}}
\begin{figure}[htbp]
\centering
\capstart

\noindent\sphinxincludegraphics[width=0.800\linewidth]{{slope_field_q2b}.png}
\caption{Slope field for \(y' = t^2 - 1\).}\label{\detokenize{Language/Science_language/ODE_calculus/course_MATH_3035/homework_assignments/index_pdf:id9}}\end{figure}




\section{Solution of Problem 3}
\label{\detokenize{Language/Science_language/ODE_calculus/course_MATH_3035/homework_assignments/index_pdf:solution-of-problem-3}}
\sphinxAtStartPar
Question: Solve the initial value problem
\begin{quote}
\begin{equation*}
\begin{split}y'=\frac{1}{t^2-3t+2},
\qquad
y(0)=1.\end{split}
\end{equation*}\end{quote}

\sphinxAtStartPar
Answer:


\subsection{Step 1: Rewrite the Differential Equation}
\label{\detokenize{Language/Science_language/ODE_calculus/course_MATH_3035/homework_assignments/index_pdf:step-1-rewrite-the-differential-equation}}
\sphinxAtStartPar
Recall that
\begin{equation*}
\begin{split}y'=\frac{dy}{dt}.\end{split}
\end{equation*}
\sphinxAtStartPar
Therefore,
\begin{equation*}
\begin{split}\frac{dy}{dt}=\frac{1}{t^2-3t+2}.\end{split}
\end{equation*}
\sphinxAtStartPar
To recover the function \(y(t)\), we integrate both sides with respect
to \(t\).
\begin{equation*}
\begin{split}y=\int \frac{1}{t^2-3t+2}\,dt + C.\end{split}
\end{equation*}
\sphinxAtStartPar
The main task is therefore to evaluate the integral
\begin{equation*}
\begin{split}\int \frac{1}{t^2-3t+2}\,dt.\end{split}
\end{equation*}

\subsection{Step 2: Factor the Denominator}
\label{\detokenize{Language/Science_language/ODE_calculus/course_MATH_3035/homework_assignments/index_pdf:step-2-factor-the-denominator}}
\sphinxAtStartPar
Before using partial fractions, factor the quadratic denominator.

\sphinxAtStartPar
We seek two numbers whose
\begin{itemize}
\item {} 
\sphinxAtStartPar
product is \(2\),

\item {} 
\sphinxAtStartPar
sum is \(-3\).

\end{itemize}

\sphinxAtStartPar
These numbers are \(-1\) and \(-2\).

\sphinxAtStartPar
Therefore,
\begin{equation*}
\begin{split}t^2-3t+2=(t-1)(t-2).\end{split}
\end{equation*}
\sphinxAtStartPar
Substituting this factorization into the integral gives
\begin{equation*}
\begin{split}y=\int \frac{1}{(t-1)(t-2)}\,dt + C.\end{split}
\end{equation*}

\subsection{Step 3: Set Up the Partial Fraction Decomposition}
\label{\detokenize{Language/Science_language/ODE_calculus/course_MATH_3035/homework_assignments/index_pdf:step-3-set-up-the-partial-fraction-decomposition}}
\sphinxAtStartPar
Because the denominator consists of two distinct linear factors,
we assume that
\begin{equation*}
\begin{split}\frac{1}{(t-1)(t-2)}
=
\frac{A}{t-1}
+
\frac{B}{t-2},\end{split}
\end{equation*}
\sphinxAtStartPar
where \(A\) and \(B\) are constants to be determined.


\subsection{Step 4: Eliminate the Denominators}
\label{\detokenize{Language/Science_language/ODE_calculus/course_MATH_3035/homework_assignments/index_pdf:step-4-eliminate-the-denominators}}
\sphinxAtStartPar
Multiply both sides by the common denominator
\begin{equation*}
\begin{split}(t-1)(t-2).\end{split}
\end{equation*}
\sphinxAtStartPar
This gives
\begin{equation*}
\begin{split}(t-1)(t-2)
\left(
\frac{1}{(t-1)(t-2)}
\right)
=
(t-1)(t-2)
\left(
\frac{A}{t-1}
+
\frac{B}{t-2}
\right).\end{split}
\end{equation*}
\sphinxAtStartPar
Simplifying each side,

\sphinxAtStartPar
Left\sphinxhyphen{}hand side:
\begin{equation*}
\begin{split}1.\end{split}
\end{equation*}
\sphinxAtStartPar
Right\sphinxhyphen{}hand side:
\begin{equation*}
\begin{split}A(t-2)+B(t-1).\end{split}
\end{equation*}
\sphinxAtStartPar
Hence,
\begin{equation*}
\begin{split}1=A(t-2)+B(t-1).\end{split}
\end{equation*}
\sphinxAtStartPar
This equation must be true for every value of \(t\).


\subsection{Step 5: Find the Constant A}
\label{\detokenize{Language/Science_language/ODE_calculus/course_MATH_3035/homework_assignments/index_pdf:step-5-find-the-constant-a}}
\sphinxAtStartPar
Choose
\begin{equation*}
\begin{split}t=1.\end{split}
\end{equation*}
\sphinxAtStartPar
Substituting into
\begin{equation*}
\begin{split}1=A(t-2)+B(t-1)\end{split}
\end{equation*}
\sphinxAtStartPar
gives
\begin{equation*}
\begin{split}1=A(1-2)+B(1-1).\end{split}
\end{equation*}
\sphinxAtStartPar
Simplify:
\begin{equation*}
\begin{split}1=A(-1)+B(0).\end{split}
\end{equation*}\begin{equation*}
\begin{split}1=-A.\end{split}
\end{equation*}
\sphinxAtStartPar
Therefore,
\begin{equation*}
\begin{split}A=-1.\end{split}
\end{equation*}

\subsection{Step 6: Find the Constant B}
\label{\detokenize{Language/Science_language/ODE_calculus/course_MATH_3035/homework_assignments/index_pdf:step-6-find-the-constant-b}}
\sphinxAtStartPar
Choose
\begin{equation*}
\begin{split}t=2.\end{split}
\end{equation*}
\sphinxAtStartPar
Substituting into
\begin{equation*}
\begin{split}1=A(t-2)+B(t-1)\end{split}
\end{equation*}
\sphinxAtStartPar
gives
\begin{equation*}
\begin{split}1=A(2-2)+B(2-1).\end{split}
\end{equation*}
\sphinxAtStartPar
Simplify:
\begin{equation*}
\begin{split}1=A(0)+B(1).\end{split}
\end{equation*}\begin{equation*}
\begin{split}1=B.\end{split}
\end{equation*}
\sphinxAtStartPar
Therefore,
\begin{equation*}
\begin{split}B=1.\end{split}
\end{equation*}

\subsection{Step 7: Write the Partial Fraction Expansion}
\label{\detokenize{Language/Science_language/ODE_calculus/course_MATH_3035/homework_assignments/index_pdf:step-7-write-the-partial-fraction-expansion}}
\sphinxAtStartPar
Substituting
\begin{equation*}
\begin{split}A=-1,
\qquad
B=1,\end{split}
\end{equation*}
\sphinxAtStartPar
into the decomposition gives
\begin{equation*}
\begin{split}\frac{1}{(t-1)(t-2)}
=
-\frac{1}{t-1}
+
\frac{1}{t-2}.\end{split}
\end{equation*}
\sphinxAtStartPar
Therefore,
\begin{equation*}
\begin{split}y
=
\int
\left(
-\frac{1}{t-1}
+
\frac{1}{t-2}
\right)
dt
+C.\end{split}
\end{equation*}

\subsection{Step 8: Integrate the First Term}
\label{\detokenize{Language/Science_language/ODE_calculus/course_MATH_3035/homework_assignments/index_pdf:step-8-integrate-the-first-term}}
\sphinxAtStartPar
Consider
\begin{equation*}
\begin{split}\int -\frac{1}{t-1}\,dt.\end{split}
\end{equation*}
\sphinxAtStartPar
Let
\begin{equation*}
\begin{split}u=t-1.\end{split}
\end{equation*}
\sphinxAtStartPar
Differentiate both sides:
\begin{equation*}
\begin{split}\frac{du}{dt}
=
\frac{d}{dt}(t-1).\end{split}
\end{equation*}
\sphinxAtStartPar
Since
\begin{equation*}
\begin{split}\frac{d}{dt}(t)=1,
\qquad
\frac{d}{dt}(-1)=0,\end{split}
\end{equation*}
\sphinxAtStartPar
we obtain
\begin{equation*}
\begin{split}\frac{du}{dt}=1.\end{split}
\end{equation*}
\sphinxAtStartPar
Multiplying by \(dt\) gives
\begin{equation*}
\begin{split}du=dt.\end{split}
\end{equation*}
\sphinxAtStartPar
Therefore,
\begin{equation*}
\begin{split}\int -\frac{1}{t-1}\,dt
=
-\int \frac{1}{u}\,du.\end{split}
\end{equation*}
\sphinxAtStartPar
Using the logarithmic integration rule
\begin{equation*}
\begin{split}\int \frac{1}{u}\,du=\ln|u|,\end{split}
\end{equation*}
\sphinxAtStartPar
we obtain
\begin{equation*}
\begin{split}-\ln|u|.\end{split}
\end{equation*}
\sphinxAtStartPar
Substituting back \(u=t-1\),
\begin{equation*}
\begin{split}-\ln|t-1|.\end{split}
\end{equation*}

\subsection{Step 9: Integrate the Second Term}
\label{\detokenize{Language/Science_language/ODE_calculus/course_MATH_3035/homework_assignments/index_pdf:step-9-integrate-the-second-term}}
\sphinxAtStartPar
Consider
\begin{equation*}
\begin{split}\int \frac{1}{t-2}\,dt.\end{split}
\end{equation*}
\sphinxAtStartPar
Let
\begin{equation*}
\begin{split}u=t-2.\end{split}
\end{equation*}
\sphinxAtStartPar
Differentiate:
\begin{equation*}
\begin{split}du=dt.\end{split}
\end{equation*}
\sphinxAtStartPar
Therefore,
\begin{equation*}
\begin{split}\int \frac{1}{t-2}\,dt
=
\int \frac{1}{u}\,du.\end{split}
\end{equation*}
\sphinxAtStartPar
Using
\begin{equation*}
\begin{split}\int \frac{1}{u}\,du=\ln|u|,\end{split}
\end{equation*}
\sphinxAtStartPar
we obtain
\begin{equation*}
\begin{split}\ln|u|.\end{split}
\end{equation*}
\sphinxAtStartPar
Substituting back \(u=t-2\),
\begin{equation*}
\begin{split}\ln|t-2|.\end{split}
\end{equation*}

\subsection{Step 10: Combine the Results}
\label{\detokenize{Language/Science_language/ODE_calculus/course_MATH_3035/homework_assignments/index_pdf:step-10-combine-the-results}}
\sphinxAtStartPar
Adding the two antiderivatives,
\begin{equation*}
\begin{split}y
=
-\ln|t-1|
+
\ln|t-2|
+
C.\end{split}
\end{equation*}

\subsection{Step 11: Combine the Logarithms}
\label{\detokenize{Language/Science_language/ODE_calculus/course_MATH_3035/homework_assignments/index_pdf:step-11-combine-the-logarithms}}
\sphinxAtStartPar
Use the logarithm rule
\begin{equation*}
\begin{split}\ln(A)-\ln(B)
=
\ln\left(\frac{A}{B}\right).\end{split}
\end{equation*}
\sphinxAtStartPar
Therefore,
\begin{align*}\!\begin{aligned}
-\ln|t-1|
+
\ln|t-2|\\
=
\ln|t-2|
-
\ln|t-1|\\
=
\ln\left|
\frac{t-2}{t-1}
\right|.\\
\end{aligned}\end{align*}
\sphinxAtStartPar
Thus,
\begin{equation*}
\begin{split}y
=
\ln\left|
\frac{t-2}{t-1}
\right|
+C.\end{split}
\end{equation*}

\subsection{Step 12: Apply the Initial Condition}
\label{\detokenize{Language/Science_language/ODE_calculus/course_MATH_3035/homework_assignments/index_pdf:step-12-apply-the-initial-condition}}
\sphinxAtStartPar
The initial condition is
\begin{equation*}
\begin{split}y(0)=1.\end{split}
\end{equation*}
\sphinxAtStartPar
Substitute \(t=0\) into the solution.
\begin{equation*}
\begin{split}1
=
\ln\left|
\frac{0-2}{0-1}
\right|
+C.\end{split}
\end{equation*}
\sphinxAtStartPar
Simplify the fraction:
\begin{equation*}
\begin{split}1
=
\ln\left|
\frac{-2}{-1}
\right|
+C.\end{split}
\end{equation*}\begin{equation*}
\begin{split}1
=
\ln(2)
+C.\end{split}
\end{equation*}
\sphinxAtStartPar
Solve for \(C\).
\begin{equation*}
\begin{split}C=1-\ln(2).\end{split}
\end{equation*}

\subsection{Step 13: Final Solution}
\label{\detokenize{Language/Science_language/ODE_calculus/course_MATH_3035/homework_assignments/index_pdf:step-13-final-solution}}
\sphinxAtStartPar
Substitute the value of \(C\) into the general solution.
\begin{equation*}
\begin{split}y(t)
=
\ln\left|
\frac{t-2}{t-1}
\right|
+1-\ln(2).\end{split}
\end{equation*}
\sphinxAtStartPar
Therefore,
\begin{equation*}
\begin{split}\boxed{
y(t)
=
\ln\left|
\frac{t-2}{t-1}
\right|
+1-\ln(2)
}.\end{split}
\end{equation*}



\section{Solution of Problem 4}
\label{\detokenize{Language/Science_language/ODE_calculus/course_MATH_3035/homework_assignments/index_pdf:solution-of-problem-4}}
\sphinxAtStartPar
\sphinxstylestrong{Given:}
\begin{equation*}
\begin{split}t^2y'' + 5ty' + 4y = 0,
\qquad
y(1)=0\end{split}
\end{equation*}
\sphinxAtStartPar
Proposed solution:
\begin{equation*}
\begin{split}y_1(t)=t^{-2}\ln t\end{split}
\end{equation*}

\subsection{Objective}
\label{\detokenize{Language/Science_language/ODE_calculus/course_MATH_3035/homework_assignments/index_pdf:objective}}
\sphinxAtStartPar
Verify that \(y_1(t)=t^{-2}\ln t\) satisfies both:
\begin{enumerate}
\sphinxsetlistlabels{\arabic}{enumi}{enumii}{}{.}%
\item {} 
\sphinxAtStartPar
The differential equation
\begin{equation*}
\begin{split}t^2y'' + 5ty' + 4y = 0\end{split}
\end{equation*}
\item {} 
\sphinxAtStartPar
The initial condition
\begin{equation*}
\begin{split}y(1)=0\end{split}
\end{equation*}
\end{enumerate}


\subsection{Solution}
\label{\detokenize{Language/Science_language/ODE_calculus/course_MATH_3035/homework_assignments/index_pdf:solution}}

\subsubsection{Step 1: Compute the First Derivative}
\label{\detokenize{Language/Science_language/ODE_calculus/course_MATH_3035/homework_assignments/index_pdf:step-1-compute-the-first-derivative}}
\sphinxAtStartPar
Given
\begin{equation*}
\begin{split}y_1=t^{-2}\ln t\end{split}
\end{equation*}
\sphinxAtStartPar
Using the product rule,
\begin{equation*}
\begin{split}\frac{d}{dt}(uv)=u'v+uv'\end{split}
\end{equation*}
\sphinxAtStartPar
Let
\begin{equation*}
\begin{split}u=t^{-2},
\qquad
v=\ln t\end{split}
\end{equation*}
\sphinxAtStartPar
Then
\begin{equation*}
\begin{split}u'=-2t^{-3},
\qquad
v'=\frac{1}{t}\end{split}
\end{equation*}
\sphinxAtStartPar
Therefore,
\begin{equation*}
\begin{split}y_1'
=
(-2t^{-3})(\ln t)
+
t^{-2}\left(\frac{1}{t}\right)\end{split}
\end{equation*}\begin{equation*}
\begin{split}y_1'
=
-2t^{-3}\ln t+t^{-3}\end{split}
\end{equation*}
\sphinxAtStartPar
Factoring out \(t^{-3}\) gives
\begin{equation*}
\begin{split}y_1'
=
t^{-3}(1-2\ln t)\end{split}
\end{equation*}

\subsubsection{Step 2: Compute the Second Derivative}
\label{\detokenize{Language/Science_language/ODE_calculus/course_MATH_3035/homework_assignments/index_pdf:step-2-compute-the-second-derivative}}
\sphinxAtStartPar
Differentiate
\begin{equation*}
\begin{split}y_1'
=
t^{-3}(1-2\ln t)\end{split}
\end{equation*}
\sphinxAtStartPar
Again using the product rule, let
\begin{equation*}
\begin{split}u=t^{-3},
\qquad
v=(1-2\ln t)\end{split}
\end{equation*}
\sphinxAtStartPar
Then
\begin{equation*}
\begin{split}u'=-3t^{-4}\end{split}
\end{equation*}
\sphinxAtStartPar
and
\begin{equation*}
\begin{split}v'=-\frac{2}{t}\end{split}
\end{equation*}
\sphinxAtStartPar
Therefore,
\begin{equation*}
\begin{split}y_1''
=
(-3t^{-4})(1-2\ln t)
+
t^{-3}\left(-\frac{2}{t}\right)\end{split}
\end{equation*}\begin{equation*}
\begin{split}y_1''
=
-3t^{-4}(1-2\ln t)-2t^{-4}\end{split}
\end{equation*}
\sphinxAtStartPar
Expanding,
\begin{equation*}
\begin{split}y_1''
=
-3t^{-4}+6t^{-4}\ln t-2t^{-4}\end{split}
\end{equation*}
\sphinxAtStartPar
Combining like terms,
\begin{equation*}
\begin{split}y_1''
=
t^{-4}(6\ln t-5)\end{split}
\end{equation*}

\subsubsection{Step 3: Substitute into the Differential Equation}
\label{\detokenize{Language/Science_language/ODE_calculus/course_MATH_3035/homework_assignments/index_pdf:step-3-substitute-into-the-differential-equation}}
\sphinxAtStartPar
The differential equation is
\begin{equation*}
\begin{split}t^2y''+5ty'+4y\end{split}
\end{equation*}
\sphinxAtStartPar
Compute each term separately.

\sphinxAtStartPar
First term:
\begin{equation*}
\begin{split}t^2y_1''
=
t^2\Big[t^{-4}(6\ln t-5)\Big]\end{split}
\end{equation*}\begin{equation*}
\begin{split}=
t^{-2}(6\ln t-5)\end{split}
\end{equation*}
\sphinxAtStartPar
Second term:
\begin{equation*}
\begin{split}5ty_1'
=
5t\Big[t^{-3}(1-2\ln t)\Big]\end{split}
\end{equation*}\begin{equation*}
\begin{split}=
5t^{-2}(1-2\ln t)\end{split}
\end{equation*}
\sphinxAtStartPar
Third term:
\begin{equation*}
\begin{split}4y_1
=
4t^{-2}\ln t\end{split}
\end{equation*}
\sphinxAtStartPar
Adding the three terms,
\begin{equation*}
\begin{split}t^2y_1''+5ty_1'+4y_1\end{split}
\end{equation*}\begin{equation*}
\begin{split}=
t^{-2}(6\ln t-5)
+
5t^{-2}(1-2\ln t)
+
4t^{-2}\ln t\end{split}
\end{equation*}
\sphinxAtStartPar
Factor out \(t^{-2}\):
\begin{equation*}
\begin{split}=
t^{-2}
\Big[
(6\ln t-5)
+5(1-2\ln t)
+4\ln t
\Big]\end{split}
\end{equation*}
\sphinxAtStartPar
Expand:
\begin{equation*}
\begin{split}=
t^{-2}
\Big[
6\ln t-5
+5
-10\ln t
+4\ln t
\Big]\end{split}
\end{equation*}
\sphinxAtStartPar
Combine like terms:
\begin{equation*}
\begin{split}=
t^{-2}
\Big[
(6-10+4)\ln t
+(-5+5)
\Big]\end{split}
\end{equation*}\begin{equation*}
\begin{split}=
t^{-2}(0)\end{split}
\end{equation*}\begin{equation*}
\begin{split}=0\end{split}
\end{equation*}
\sphinxAtStartPar
Hence,
\begin{equation*}
\begin{split}t^2y_1''+5ty_1'+4y_1=0\end{split}
\end{equation*}
\sphinxAtStartPar
Therefore, \(y_1\) satisfies the differential equation.


\subsubsection{Step 4: Verify the Initial Condition}
\label{\detokenize{Language/Science_language/ODE_calculus/course_MATH_3035/homework_assignments/index_pdf:step-4-verify-the-initial-condition}}
\sphinxAtStartPar
Evaluate \(y_1(1)\):
\begin{equation*}
\begin{split}y_1(1)
=
1^{-2}\ln(1)\end{split}
\end{equation*}
\sphinxAtStartPar
Since
\begin{equation*}
\begin{split}\ln(1)=0\end{split}
\end{equation*}
\sphinxAtStartPar
it follows that
\begin{equation*}
\begin{split}y_1(1)=0\end{split}
\end{equation*}
\sphinxAtStartPar
Thus the initial condition is satisfied.


\subsubsection{Conclusion}
\label{\detokenize{Language/Science_language/ODE_calculus/course_MATH_3035/homework_assignments/index_pdf:conclusion}}
\sphinxAtStartPar
Since
\begin{equation*}
\begin{split}t^2y_1''+5ty_1'+4y_1=0\end{split}
\end{equation*}
\sphinxAtStartPar
and
\begin{equation*}
\begin{split}y_1(1)=0,\end{split}
\end{equation*}
\sphinxAtStartPar
the function
\begin{equation*}
\begin{split}y_1(t)=t^{-2}\ln t\end{split}
\end{equation*}
\sphinxAtStartPar
satisfies both the differential equation and the initial condition.

\sphinxAtStartPar
Therefore,
\begin{equation*}
\begin{split}\boxed{y_1(t)=t^{-2}\ln t}\end{split}
\end{equation*}
\sphinxAtStartPar
is a solution of the given initial value problem.




\section{Solution of Problem 5}
\label{\detokenize{Language/Science_language/ODE_calculus/course_MATH_3035/homework_assignments/index_pdf:solution-of-problem-5}}

\subsection{Problem}
\label{\detokenize{Language/Science_language/ODE_calculus/course_MATH_3035/homework_assignments/index_pdf:problem}}
\sphinxAtStartPar
Solve the differential equation
\begin{equation*}
\begin{split}y' + 2y = \frac{e^{-2t}}{1+t^2}.\end{split}
\end{equation*}

\subsection{Solution}
\label{\detokenize{Language/Science_language/ODE_calculus/course_MATH_3035/homework_assignments/index_pdf:id2}}
\sphinxAtStartPar
This is a first\sphinxhyphen{}order linear differential equation.

\sphinxAtStartPar
Recall the standard linear form:
\begin{equation*}
\begin{split}y' + P(t)y = Q(t).\end{split}
\end{equation*}
\sphinxAtStartPar
Comparing with the given equation,
\begin{equation*}
\begin{split}y' + 2y = \frac{e^{-2t}}{1+t^2},\end{split}
\end{equation*}
\sphinxAtStartPar
we identify
\begin{equation*}
\begin{split}P(t)=2,\end{split}
\end{equation*}
\sphinxAtStartPar
and
\begin{equation*}
\begin{split}Q(t)=\frac{e^{-2t}}{1+t^2}.\end{split}
\end{equation*}

\bigskip\hrule\bigskip



\subsubsection{Step 1: Find the Integrating Factor}
\label{\detokenize{Language/Science_language/ODE_calculus/course_MATH_3035/homework_assignments/index_pdf:step-1-find-the-integrating-factor}}
\sphinxAtStartPar
The integrating factor is
\begin{equation*}
\begin{split}\mu(t)=e^{\int P(t)\,dt}.\end{split}
\end{equation*}
\sphinxAtStartPar
Substituting \(P(t)=2\),
\begin{equation*}
\begin{split}\mu(t)=e^{\int 2\,dt}.\end{split}
\end{equation*}
\sphinxAtStartPar
Evaluating the integral,
\begin{equation*}
\begin{split}\mu(t)=e^{2t}.\end{split}
\end{equation*}
\sphinxAtStartPar
Therefore, the integrating factor is
\begin{equation*}
\begin{split}\boxed{\mu(t)=e^{2t}}.\end{split}
\end{equation*}

\bigskip\hrule\bigskip



\subsubsection{Step 2: Multiply the Differential Equation by the Integrating Factor}
\label{\detokenize{Language/Science_language/ODE_calculus/course_MATH_3035/homework_assignments/index_pdf:step-2-multiply-the-differential-equation-by-the-integrating-factor}}
\sphinxAtStartPar
Multiply every term of the differential equation by \(e^{2t}\):
\begin{equation*}
\begin{split}e^{2t}y' + 2e^{2t}y
=
e^{2t}\left(\frac{e^{-2t}}{1+t^2}\right).\end{split}
\end{equation*}
\sphinxAtStartPar
Simplifying the right\sphinxhyphen{}hand side,
\begin{equation*}
\begin{split}e^{2t}e^{-2t}=1.\end{split}
\end{equation*}
\sphinxAtStartPar
Therefore,
\begin{equation*}
\begin{split}e^{2t}y' + 2e^{2t}y
=
\frac{1}{1+t^2}.\end{split}
\end{equation*}

\bigskip\hrule\bigskip



\subsubsection{Step 3: Recognize the Product Derivative}
\label{\detokenize{Language/Science_language/ODE_calculus/course_MATH_3035/homework_assignments/index_pdf:step-3-recognize-the-product-derivative}}
\sphinxAtStartPar
Recall the product rule:
\begin{equation*}
\begin{split}\frac{d}{dt}(uv)
=
u\frac{dv}{dt}
+
v\frac{du}{dt}.\end{split}
\end{equation*}
\sphinxAtStartPar
Let
\begin{equation*}
\begin{split}u=e^{2t},
\qquad
v=y.\end{split}
\end{equation*}
\sphinxAtStartPar
Then
\begin{equation*}
\begin{split}\frac{du}{dt}=2e^{2t}.\end{split}
\end{equation*}
\sphinxAtStartPar
Applying the product rule,
\begin{equation*}
\begin{split}\frac{d}{dt}(e^{2t}y)
=
e^{2t}y'
+
2e^{2t}y.\end{split}
\end{equation*}
\sphinxAtStartPar
This matches the left\sphinxhyphen{}hand side of the differential equation exactly.

\sphinxAtStartPar
Hence,
\begin{equation*}
\begin{split}\frac{d}{dt}(e^{2t}y)
=
\frac{1}{1+t^2}.\end{split}
\end{equation*}

\bigskip\hrule\bigskip



\subsubsection{Step 4: Integrate Both Sides}
\label{\detokenize{Language/Science_language/ODE_calculus/course_MATH_3035/homework_assignments/index_pdf:step-4-integrate-both-sides}}
\sphinxAtStartPar
Integrating both sides with respect to \(t\),
\begin{equation*}
\begin{split}\int \frac{d}{dt}(e^{2t}y)\,dt
=
\int \frac{1}{1+t^2}\,dt.\end{split}
\end{equation*}
\sphinxAtStartPar
The left\sphinxhyphen{}hand side simplifies to
\begin{equation*}
\begin{split}e^{2t}y.\end{split}
\end{equation*}
\sphinxAtStartPar
Using the standard antiderivative
\begin{equation*}
\begin{split}\int \frac{1}{1+t^2}\,dt
=
\tan^{-1}(t)+C,\end{split}
\end{equation*}
\sphinxAtStartPar
we obtain
\begin{equation*}
\begin{split}e^{2t}y
=
\tan^{-1}(t)+C.\end{split}
\end{equation*}

\bigskip\hrule\bigskip



\subsubsection{Step 5: Solve for \protect\(y\protect\)}
\label{\detokenize{Language/Science_language/ODE_calculus/course_MATH_3035/homework_assignments/index_pdf:step-5-solve-for-y}}
\sphinxAtStartPar
Multiply both sides by \(e^{-2t}\):
\begin{equation*}
\begin{split}y
=
e^{-2t}\bigl(\tan^{-1}(t)+C\bigr).\end{split}
\end{equation*}
\sphinxAtStartPar
Therefore,
\begin{equation*}
\begin{split}\boxed{
y(t)
=
e^{-2t}\left(\tan^{-1}(t)+C\right)
}.\end{split}
\end{equation*}

\subsection{Answer}
\label{\detokenize{Language/Science_language/ODE_calculus/course_MATH_3035/homework_assignments/index_pdf:answer}}\begin{equation*}
\begin{split}\boxed{
y(t)
=
e^{-2t}\left(\tan^{-1}(t)+C\right)
}.\end{split}
\end{equation*}



\section{Solution of Problem 6}
\label{\detokenize{Language/Science_language/ODE_calculus/course_MATH_3035/homework_assignments/index_pdf:solution-of-problem-6}}

\subsection{Problem}
\label{\detokenize{Language/Science_language/ODE_calculus/course_MATH_3035/homework_assignments/index_pdf:id3}}
\sphinxAtStartPar
Solve the differential equation
\begin{equation*}
\begin{split}y' - \frac{1}{t}\sin(t^2) = -\frac{2}{t}y.\end{split}
\end{equation*}

\bigskip\hrule\bigskip



\subsection{Step 1: Put the Equation into Standard Linear Form}
\label{\detokenize{Language/Science_language/ODE_calculus/course_MATH_3035/homework_assignments/index_pdf:step-1-put-the-equation-into-standard-linear-form}}
\sphinxAtStartPar
Add \(\frac{2}{t}y\) to both sides:
\begin{equation*}
\begin{split}y' + \frac{2}{t}y = \frac{\sin(t^2)}{t}.\end{split}
\end{equation*}
\sphinxAtStartPar
Compare with the standard linear form
\begin{equation*}
\begin{split}y' + P(t)y = Q(t).\end{split}
\end{equation*}
\sphinxAtStartPar
Thus,
\begin{equation*}
\begin{split}P(t)=\frac{2}{t},\end{split}
\end{equation*}
\sphinxAtStartPar
and
\begin{equation*}
\begin{split}Q(t)=\frac{\sin(t^2)}{t}.\end{split}
\end{equation*}

\bigskip\hrule\bigskip



\subsection{Step 2: Find the Integrating Factor}
\label{\detokenize{Language/Science_language/ODE_calculus/course_MATH_3035/homework_assignments/index_pdf:step-2-find-the-integrating-factor}}
\sphinxAtStartPar
The integrating factor is
\begin{equation*}
\begin{split}\mu(t)=e^{\int P(t)\,dt}.\end{split}
\end{equation*}
\sphinxAtStartPar
Substituting \(P(t)=\frac{2}{t}\) gives
\begin{equation*}
\begin{split}\mu(t)
=
e^{\int \frac{2}{t}\,dt}.\end{split}
\end{equation*}
\sphinxAtStartPar
Using
\begin{equation*}
\begin{split}\int \frac{1}{t}\,dt = \ln|t|,\end{split}
\end{equation*}
\sphinxAtStartPar
we obtain
\begin{equation*}
\begin{split}\mu(t)
=
e^{2\ln|t|}.\end{split}
\end{equation*}
\sphinxAtStartPar
Using the identity
\begin{equation*}
\begin{split}e^{\ln a}=a,\end{split}
\end{equation*}
\sphinxAtStartPar
we get
\begin{equation*}
\begin{split}\mu(t)=|t|^2.\end{split}
\end{equation*}
\sphinxAtStartPar
For \(t>0\),
\begin{equation*}
\begin{split}\boxed{\mu(t)=t^2}.\end{split}
\end{equation*}

\bigskip\hrule\bigskip



\subsection{Step 3: Multiply the Equation by the Integrating Factor}
\label{\detokenize{Language/Science_language/ODE_calculus/course_MATH_3035/homework_assignments/index_pdf:step-3-multiply-the-equation-by-the-integrating-factor}}
\sphinxAtStartPar
Multiply every term by \(t^2\):
\begin{equation*}
\begin{split}t^2y' + 2ty = t\sin(t^2).\end{split}
\end{equation*}
\sphinxAtStartPar
Notice that
\begin{equation*}
\begin{split}\frac{d}{dt}(t^2y)
=
t^2y' + 2ty.\end{split}
\end{equation*}
\sphinxAtStartPar
Therefore,
\begin{equation*}
\begin{split}\frac{d}{dt}(t^2y)
=
t\sin(t^2).\end{split}
\end{equation*}

\bigskip\hrule\bigskip



\subsection{Step 4: Integrate Both Sides}
\label{\detokenize{Language/Science_language/ODE_calculus/course_MATH_3035/homework_assignments/index_pdf:id4}}
\sphinxAtStartPar
Integrating both sides gives
\begin{equation*}
\begin{split}t^2y
=
\int t\sin(t^2)\,dt + C.\end{split}
\end{equation*}

\bigskip\hrule\bigskip



\subsubsection{Step 4a: Use a Substitution}
\label{\detokenize{Language/Science_language/ODE_calculus/course_MATH_3035/homework_assignments/index_pdf:step-4a-use-a-substitution}}
\sphinxAtStartPar
Let
\begin{equation*}
\begin{split}u=t^2.\end{split}
\end{equation*}
\sphinxAtStartPar
Differentiate:
\begin{equation*}
\begin{split}du = 2t\,dt.\end{split}
\end{equation*}
\sphinxAtStartPar
Solve for \(t\,dt\):
\begin{equation*}
\begin{split}t\,dt=\frac{du}{2}.\end{split}
\end{equation*}
\sphinxAtStartPar
Substitute into the integral:
\begin{equation*}
\begin{split}\int t\sin(t^2)\,dt = \frac12\int \sin(u)\,du.\end{split}
\end{equation*}

\bigskip\hrule\bigskip



\subsubsection{Step 4b: Evaluate the Integral}
\label{\detokenize{Language/Science_language/ODE_calculus/course_MATH_3035/homework_assignments/index_pdf:step-4b-evaluate-the-integral}}
\sphinxAtStartPar
Recall the basic antiderivative
\begin{equation*}
\begin{split}\int \sin(u)\,du
=
-\cos(u)+C.\end{split}
\end{equation*}
\sphinxAtStartPar
Therefore,
\begin{equation*}
\begin{split}\frac12\int \sin(u)\,du = -\frac12\cos(u).\end{split}
\end{equation*}
\sphinxAtStartPar
Substituting back \(u=t^2\) gives
\begin{equation*}
\begin{split}\int t\sin(t^2)\,dt = -\frac12\cos(t^2).\end{split}
\end{equation*}
\sphinxAtStartPar
Hence
\begin{equation*}
\begin{split}t^2y = -\frac12\cos(t^2)+C.\end{split}
\end{equation*}

\bigskip\hrule\bigskip



\subsection{Step 5: Solve for \protect\(y\protect\)}
\label{\detokenize{Language/Science_language/ODE_calculus/course_MATH_3035/homework_assignments/index_pdf:id5}}
\sphinxAtStartPar
Divide both sides by \(t^2\):
\begin{align*}\!\begin{aligned}
y(t)\\
=\\
\frac{C-\frac12\cos(t^2)}{t^2}.\\
\end{aligned}\end{align*}
\sphinxAtStartPar
Equivalent form:
\begin{equation*}
\begin{split}\boxed{
y(t)
=
\frac{C}{t^2}
-
\frac{\cos(t^2)}{2t^2}
}.\end{split}
\end{equation*}



\section{Solution of Problem 7}
\label{\detokenize{Language/Science_language/ODE_calculus/course_MATH_3035/homework_assignments/index_pdf:solution-of-problem-7}}

\subsection{Problem}
\label{\detokenize{Language/Science_language/ODE_calculus/course_MATH_3035/homework_assignments/index_pdf:id6}}
\sphinxAtStartPar
Solve the differential equation
\begin{equation*}
\begin{split}y' + \cot(t)\,y = e^t\cot(t).\end{split}
\end{equation*}

\subsection{Step 1: Verify That the Equation Is Linear}
\label{\detokenize{Language/Science_language/ODE_calculus/course_MATH_3035/homework_assignments/index_pdf:step-1-verify-that-the-equation-is-linear}}
\sphinxAtStartPar
A first\sphinxhyphen{}order linear differential equation has the form
\begin{equation*}
\begin{split}y' + P(t)y = Q(t).\end{split}
\end{equation*}
\sphinxAtStartPar
Comparing with
\begin{equation*}
\begin{split}y' + \cot(t)\,y = e^t\cot(t),\end{split}
\end{equation*}
\sphinxAtStartPar
we identify
\begin{equation*}
\begin{split}P(t)=\cot(t)\end{split}
\end{equation*}
\sphinxAtStartPar
and
\begin{equation*}
\begin{split}Q(t)=e^t\cot(t).\end{split}
\end{equation*}
\sphinxAtStartPar
Since the equation is in standard linear form, we can use the integrating factor method.


\bigskip\hrule\bigskip



\subsection{Step 2: Find the Integrating Factor}
\label{\detokenize{Language/Science_language/ODE_calculus/course_MATH_3035/homework_assignments/index_pdf:id7}}
\sphinxAtStartPar
The integrating factor is defined by
\begin{equation*}
\begin{split}\mu(t)=e^{\int P(t)\,dt}.\end{split}
\end{equation*}
\sphinxAtStartPar
Substituting \(P(t)=\cot(t)\) gives
\begin{equation*}
\begin{split}\mu(t)=e^{\int \cot(t)\,dt}.\end{split}
\end{equation*}
\sphinxAtStartPar
Recall that
\begin{equation*}
\begin{split}\cot(t)=\frac{\cos(t)}{\sin(t)}.\end{split}
\end{equation*}
\sphinxAtStartPar
Therefore,
\begin{equation*}
\begin{split}\int \cot(t)\,dt
=
\int \frac{\cos(t)}{\sin(t)}\,dt.\end{split}
\end{equation*}
\sphinxAtStartPar
Let
\begin{equation*}
\begin{split}u=\sin(t).\end{split}
\end{equation*}
\sphinxAtStartPar
Then
\begin{equation*}
\begin{split}du=\cos(t)\,dt.\end{split}
\end{equation*}
\sphinxAtStartPar
Substituting,
\begin{equation*}
\begin{split}\int \frac{\cos(t)}{\sin(t)}\,dt
=
\int \frac{1}{u}\,du.\end{split}
\end{equation*}
\sphinxAtStartPar
Using the logarithm rule,
\begin{equation*}
\begin{split}\int \frac{1}{u}\,du
=
\ln|u|+C.\end{split}
\end{equation*}
\sphinxAtStartPar
Replacing \(u\) with \(\sin(t)\),
\begin{equation*}
\begin{split}\int \cot(t)\,dt
=
\ln|\sin(t)|.\end{split}
\end{equation*}
\sphinxAtStartPar
Thus,
\begin{equation*}
\begin{split}\mu(t)
=
e^{\ln|\sin(t)|}.\end{split}
\end{equation*}
\sphinxAtStartPar
Using the property
\begin{equation*}
\begin{split}e^{\ln(x)}=x,\end{split}
\end{equation*}
\sphinxAtStartPar
we obtain
\begin{equation*}
\begin{split}\mu(t)=|\sin(t)|.\end{split}
\end{equation*}
\sphinxAtStartPar
On an interval where \(\sin(t)>0\), we take
\begin{equation*}
\begin{split}\boxed{\mu(t)=\sin(t)}.\end{split}
\end{equation*}

\bigskip\hrule\bigskip



\subsection{Step 3: Multiply the Entire Equation by the Integrating Factor}
\label{\detokenize{Language/Science_language/ODE_calculus/course_MATH_3035/homework_assignments/index_pdf:step-3-multiply-the-entire-equation-by-the-integrating-factor}}
\sphinxAtStartPar
Multiply each term by \(\sin(t)\):
\begin{equation*}
\begin{split}\sin(t)y'
+
\sin(t)\cot(t)y
=
e^t\sin(t)\cot(t).\end{split}
\end{equation*}
\sphinxAtStartPar
Using
\begin{equation*}
\begin{split}\cot(t)=\frac{\cos(t)}{\sin(t)},\end{split}
\end{equation*}
\sphinxAtStartPar
we find
\begin{equation*}
\begin{split}\sin(t)\cot(t)
=
\sin(t)\frac{\cos(t)}{\sin(t)}
=
\cos(t).\end{split}
\end{equation*}
\sphinxAtStartPar
Therefore,
\begin{equation*}
\begin{split}\sin(t)y'
+
\cos(t)y
=
e^t\cos(t).\end{split}
\end{equation*}

\bigskip\hrule\bigskip



\subsection{Step 4: Recognize the Product Rule}
\label{\detokenize{Language/Science_language/ODE_calculus/course_MATH_3035/homework_assignments/index_pdf:step-4-recognize-the-product-rule}}
\sphinxAtStartPar
Recall the product rule:
\begin{equation*}
\begin{split}\frac{d}{dt}[u(t)v(t)]
=
u'(t)v(t)
+
u(t)v'(t).\end{split}
\end{equation*}
\sphinxAtStartPar
Let
\begin{equation*}
\begin{split}u(t)=y\end{split}
\end{equation*}
\sphinxAtStartPar
and
\begin{equation*}
\begin{split}v(t)=\sin(t).\end{split}
\end{equation*}
\sphinxAtStartPar
Then
\begin{equation*}
\begin{split}\frac{d}{dt}[y\sin(t)]
=
y'\sin(t)
+
y\cos(t).\end{split}
\end{equation*}
\sphinxAtStartPar
Notice that the left\sphinxhyphen{}hand side of our differential equation is exactly
\begin{equation*}
\begin{split}\sin(t)y'
+
\cos(t)y.\end{split}
\end{equation*}
\sphinxAtStartPar
Therefore,
\begin{equation*}
\begin{split}\frac{d}{dt}[y\sin(t)]
=
e^t\cos(t).\end{split}
\end{equation*}

\bigskip\hrule\bigskip



\subsection{Step 5: Integrate Both Sides}
\label{\detokenize{Language/Science_language/ODE_calculus/course_MATH_3035/homework_assignments/index_pdf:step-5-integrate-both-sides}}
\sphinxAtStartPar
Integrate both sides:
\begin{equation*}
\begin{split}\int \frac{d}{dt}[y\sin(t)]\,dt
=
\int e^t\cos(t)\,dt.\end{split}
\end{equation*}
\sphinxAtStartPar
The left\sphinxhyphen{}hand side simplifies to
\begin{equation*}
\begin{split}y\sin(t).\end{split}
\end{equation*}
\sphinxAtStartPar
Thus,
\begin{equation*}
\begin{split}y\sin(t)
=
\int e^t\cos(t)\,dt + C.\end{split}
\end{equation*}
\sphinxAtStartPar
We now evaluate the integral
\begin{equation*}
\begin{split}I=\int e^t\cos(t)\,dt.\end{split}
\end{equation*}

\bigskip\hrule\bigskip



\subsection{Step 6: Evaluate the Integral \protect\(\int e^t\cos(t)\,dt\protect\)}
\label{\detokenize{Language/Science_language/ODE_calculus/course_MATH_3035/homework_assignments/index_pdf:step-6-evaluate-the-integral-int-e-t-cos-t-dt}}

\subsubsection{Step 6a:}
\label{\detokenize{Language/Science_language/ODE_calculus/course_MATH_3035/homework_assignments/index_pdf:step-6a}}
\sphinxAtStartPar
Let
\begin{equation*}
\begin{split}I=\int e^t\cos(t)\,dt.\end{split}
\end{equation*}
\sphinxAtStartPar
Using integration by parts,
\begin{equation*}
\begin{split}\int u\,dv = uv-\int v\,du.\end{split}
\end{equation*}
\sphinxAtStartPar
Choose
\begin{equation*}
\begin{split}u=\cos(t),
\qquad
dv=e^t\,dt.\end{split}
\end{equation*}
\sphinxAtStartPar
Then
\begin{equation*}
\begin{split}du=-\sin(t)\,dt,
\qquad
v=e^t.\end{split}
\end{equation*}
\sphinxAtStartPar
Substituting,
\begin{equation*}
\begin{split}I=e^t\cos(t)-\int e^t(-\sin(t))\,dt.\end{split}
\end{equation*}
\sphinxAtStartPar
Therefore,
\begin{equation*}
\begin{split}I=e^t\cos(t)+\int e^t\sin(t)\,dt.\end{split}
\end{equation*}
\sphinxAtStartPar
Let
\begin{equation*}
\begin{split}J=\int e^t\sin(t)\,dt.\end{split}
\end{equation*}
\sphinxAtStartPar
Then
\begin{equation*}
\begin{split}I=e^t\cos(t)+J.\end{split}
\end{equation*}

\bigskip\hrule\bigskip



\subsubsection{Step 6b: Evaluate \protect\(J\protect\)}
\label{\detokenize{Language/Science_language/ODE_calculus/course_MATH_3035/homework_assignments/index_pdf:step-6b-evaluate-j}}
\sphinxAtStartPar
Again use integration by parts.

\sphinxAtStartPar
Choose
\begin{equation*}
\begin{split}u=\sin(t),
\qquad
dv=e^t\,dt.\end{split}
\end{equation*}
\sphinxAtStartPar
Then
\begin{equation*}
\begin{split}du=\cos(t)\,dt,
\qquad
v=e^t.\end{split}
\end{equation*}
\sphinxAtStartPar
Substituting,
\begin{equation*}
\begin{split}J=e^t\sin(t)-\int e^t\cos(t)\,dt.\end{split}
\end{equation*}
\sphinxAtStartPar
But
\begin{equation*}
\begin{split}\int e^t\cos(t)\,dt=I.\end{split}
\end{equation*}
\sphinxAtStartPar
Hence
\begin{equation*}
\begin{split}J=e^t\sin(t)-I.\end{split}
\end{equation*}

\bigskip\hrule\bigskip



\subsubsection{Step 6c: Substitute Back Into \protect\(I\protect\)}
\label{\detokenize{Language/Science_language/ODE_calculus/course_MATH_3035/homework_assignments/index_pdf:step-6c-substitute-back-into-i}}
\sphinxAtStartPar
Recall
\begin{equation*}
\begin{split}I=e^t\cos(t)+J.\end{split}
\end{equation*}
\sphinxAtStartPar
Substitute \(J=e^t\sin(t)-I\):
\begin{equation*}
\begin{split}I=e^t\cos(t)+e^t\sin(t)-I.\end{split}
\end{equation*}
\sphinxAtStartPar
Combine terms:
\begin{equation*}
\begin{split}I=e^t(\cos(t)+\sin(t))-I.\end{split}
\end{equation*}
\sphinxAtStartPar
Add \(I\) to both sides:
\begin{equation*}
\begin{split}2I=e^t(\cos(t)+\sin(t)).\end{split}
\end{equation*}
\sphinxAtStartPar
Divide by 2:
\begin{equation*}
\begin{split}I=\frac{e^t}{2}(\cos(t)+\sin(t))+C.\end{split}
\end{equation*}
\sphinxAtStartPar
Therefore,
\begin{equation*}
\begin{split}\boxed{
\int e^t\cos(t)\,dt
=
\frac{e^t}{2}
(\cos(t)+\sin(t))
+C
}.\end{split}
\end{equation*}

\bigskip\hrule\bigskip



\subsection{Step 7: Substitute Back Into the Solution}
\label{\detokenize{Language/Science_language/ODE_calculus/course_MATH_3035/homework_assignments/index_pdf:step-7-substitute-back-into-the-solution}}
\sphinxAtStartPar
We had
\begin{equation*}
\begin{split}y\sin(t)
=
\int e^t\cos(t)\,dt + C.\end{split}
\end{equation*}
\sphinxAtStartPar
Substituting the result,
\begin{equation*}
\begin{split}y\sin(t)
=
\frac{e^t}{2}
(\cos(t)+\sin(t))
+C.\end{split}
\end{equation*}

\bigskip\hrule\bigskip



\subsection{Step 8: Solve for \protect\(y\protect\)}
\label{\detokenize{Language/Science_language/ODE_calculus/course_MATH_3035/homework_assignments/index_pdf:step-8-solve-for-y}}
\sphinxAtStartPar
Divide both sides by \(\sin(t)\):
\begin{equation*}
\begin{split}y
=
\frac{\frac{e^t}{2}(\cos(t)+\sin(t))}
     {\sin(t)}
+
\frac{C}{\sin(t)}.\end{split}
\end{equation*}
\sphinxAtStartPar
Separate the fraction:
\begin{equation*}
\begin{split}y
=
\frac{e^t}{2}
\left(
\frac{\cos(t)}{\sin(t)}
+
\frac{\sin(t)}{\sin(t)}
\right)
+
C\csc(t).\end{split}
\end{equation*}
\sphinxAtStartPar
Since
\begin{equation*}
\begin{split}\frac{\cos(t)}{\sin(t)}
=
\cot(t)\end{split}
\end{equation*}
\sphinxAtStartPar
and
\begin{equation*}
\begin{split}\frac{\sin(t)}{\sin(t)}
=
1,\end{split}
\end{equation*}
\sphinxAtStartPar
we obtain
\begin{equation*}
\begin{split}y
=
\frac{e^t}{2}
\left(
\cot(t)+1
\right)
+
C\csc(t).\end{split}
\end{equation*}

\bigskip\hrule\bigskip



\subsection{Final Answer}
\label{\detokenize{Language/Science_language/ODE_calculus/course_MATH_3035/homework_assignments/index_pdf:final-answer}}\begin{equation*}
\begin{split}\boxed{
y(t)
=
\frac{e^t}{2}
\left(
1+\cot(t)
\right)
+
C\,\csc(t)
}\end{split}
\end{equation*}


\renewcommand{\indexname}{Index}
\printindex
\end{document}